% !TEX root = ../main.tex
%% Corrigé du devoir surveillé n° 1 — valeurs vérifiées : verifications/devoirs.py

\corrigeDevoir{1}
{\otLight\fontsize{8.5}{11.5}\selectfont\color{otMuted}\noindent Les barèmes répartissent les 20 points du devoir entre les
questions. Une réponse juste mais non justifiée, ou un résultat sans unité, ne rapporte qu'une partie des points.\par}

\begin{corrigeBac}{8}{Le bleu d'un sirop de menthe}
\begin{questions}
\item La solution absorbe surtout l'orange et le rouge (maximum vers \qty{638}{nm}) ; elle transmet le reste de la lumière blanche
et prend la \res{\text{couleur complémentaire}} du rouge-orangé : elle paraît \res{\text{bleue}}. \pts{1 pt}
\item À \qty{638}{nm}, l'absorbance est \res{\text{maximale}} : la mesure est la plus sensible et la moins affectée par une petite
erreur sur la longueur d'onde. \pts{1 pt}
\item Les points sont alignés sur une \res{\text{droite passant par l'origine}} : $A$ est proportionnelle à $c$, la loi de Beer-Lambert
est vérifiée. Coefficient directeur : $k=\dfrac{0{,}81}{5{,}0}\approx\res{\qty{0,161}{L.mg^{-1}}}$. \pts{1,5 pt}
\item Facteur 20 : $\dfrac{V_\text{fille}}{V_\text{mère}}=\dfrac{100{,}0}{5{,}0}$. On prélève \qty{5,0}{mL} de sirop avec la \res{\text{pipette jaugée de 5,0 mL}},
on les verse dans la \res{\text{fiole jaugée de 100,0 mL}}, on complète aux trois quarts avec de l'eau distillée, on agite, puis on
ajuste au trait de jauge et on homogénéise. \pts{1,5 pt}
\item Par lecture graphique ou par le calcul : $c_\text{dil}=\dfrac{A}{k}=\dfrac{0{,}52}{0{,}161}\approx\res{\qty{3,2}{mg.L^{-1}}}$ ; dans le sirop :
$c=20\times3{,}2=\res{\qty{64}{mg.L^{-1}}}$. \pts{1,5 pt}
\item Boisson : $\dfrac{64}{8}=\qty{8,0}{mg.L^{-1}}$, soit $8{,}0\times0{,}25=\qty{2,0}{mg}$ par verre. DJA : $5\times30=\qty{150}{mg}$ par jour, soit
$\dfrac{150}{2{,}0}=\res{75}$ verres : une consommation normale de sirop \res{\text{ne présente aucun risque}}. \pts{1,5 pt}
\end{questions}
\schema[corrige]{ds1-etalonnage}
\end{corrigeBac}

\begin{corrigeBac}{7}{Argenter un fil de cuivre}
\begin{questions}
\item $\ce{Ag+ + e- <=> Ag}$ et $\ce{Cu <=> Cu^{2+} + 2e-}$. L'oxydant qui réagit est l'\res{\text{ion argent}} $\ce{Ag+}$, le réducteur
le \res{\text{cuivre}} $\ce{Cu}$. \pts{1,5 pt}
\item En multipliant la première demi-équation par 2 : $\res{\ce{Cu + 2 Ag+ -> Cu^{2+} + 2 Ag}}$. \pts{0,5 pt}
\item $n(\ce{Cu})=\dfrac{0{,}635}{63{,}5}=\res{\qty{1,00e-2}{mol}}$ ; $n(\ce{Ag+})=c\,V=0{,}100\times0{,}0500=\res{\qty{5,00e-3}{mol}}$. \pts{1 pt}
\item Quantités en mmol :
\begin{avancement}{4}
\tavEquation{\ce{Cu} & $+\;2\,\ce{Ag+}$ & $\longrightarrow\;\ce{Cu^{2+}}$ & $+\;2\,\ce{Ag}$}
\tavTitres{4}
initial & $0$ & $10{,}0$ & $5{,}00$ & $0$ & $0$ \\
en cours & $x$ & $10{,}0-x$ & $5{,}00-2x$ & $x$ & $2x$ \\
final & $x_\mathrm{max}=2{,}50$ & $7{,}50$ & $0$ & $2{,}50$ & $5{,}00$ \\
\end{avancement}
Si $\ce{Cu}$ est limitant, $x=\qty{10,0}{mmol}$ ; si $\ce{Ag+}$ l'est, $5{,}00-2x=0$ donne $x=\qty{2,50}{mmol}$. On retient la plus petite valeur :
$\res{x_\mathrm{max}=\qty{2,50e-3}{mol}}$, l'\res{\text{ion argent}} est le réactif limitant. \pts{2 pts}
\item $m(\ce{Cu})=7{,}50\times10^{-3}\times63{,}5=\res{\qty{0,476}{g}}$ ; $m(\ce{Ag})=5{,}00\times10^{-3}\times107{,}9=\res{\qty{0,540}{g}}$ ;
$[\ce{Cu^{2+}}]=\dfrac{2{,}50\times10^{-3}}{0{,}0500}=\res{\qty{5,00e-2}{mol.L^{-1}}}$, d'où la teinte bleue. \pts{1,5 pt}
\item Le fil porte le cuivre restant et l'argent déposé : $0{,}476+0{,}540=\res{\qty{1,016}{g}}$. Il est plus lourd qu'au départ. \pts{0,5 pt}
\end{questions}
\retiens{Le réactif limitant se trouve en cherchant, pour chaque réactif, la valeur de $x$ qui l'épuise : la plus petite est $x_\mathrm{max}$.}
\end{corrigeBac}

\begin{corrigeBac}{5}{Une boisson sucrée}
\begin{questions}
\item $M=12\times12{,}0+22\times1{,}0+11\times16{,}0=\res{\qty{342,0}{g.mol^{-1}}}$. \pts{0,5 pt}
\item $c_m=\dfrac{m}{V}=\dfrac{35}{0{,}33}=\res{\qty{106}{g.L^{-1}}}$ ; $c=\dfrac{c_m}{M}=\dfrac{106}{342{,}0}=\res{\qty{0,310}{mol.L^{-1}}}$. \pts{1,5 pt}
\item $n=\dfrac{35}{342{,}0}=\qty{0,102}{mol}$, soit $N=n\,N_\mathrm{A}=0{,}102\times\num{6,02e23}=\res{\num{6,16e22}}$ molécules. \pts{1 pt}
\item Masse à peser : $m=c\,V\,M=0{,}310\times0{,}1000\times342{,}0=\res{\qty{10,6}{g}}$. On la pèse dans une coupelle sur une balance tarée, on
l'introduit dans une \res{\text{fiole jaugée de 100,0 mL}} avec un entonnoir, on rince coupelle et entonnoir, on dissout, puis on
complète au trait de jauge et on homogénéise. \pts{1 pt}
\item Lors d'une dilution, la quantité de soluté se conserve : $V_\text{mère}=\dfrac{c_\text{fille}\,V_\text{fille}}{c_\text{mère}}=\dfrac{50{,}0}{10}=\res{\qty{5,0}{mL}}$, prélevés
à la pipette jaugée. \pts{1 pt}
\end{questions}
\end{corrigeBac}
